\documentclass{fastverk}
\graphicspath{{../icons/}{../wordmark/}}

\begin{document}

% ── cover ───────────────────────────────────────────────────────────
\begin{center}
  \vspace*{18mm}
  \includegraphics[width=128mm]{lockup_light.png}\\[10mm]
  {\large Brand Guidelines}\\[1mm]
  {\footnotesize Generated from one parametric source — \texttt{//gen:gen\_mark}}
\end{center}
\vspace{8mm}

% ── the mark ─────────────────────────────────────────────────────────
\section*{The mark}
The mark plays off \textbf{F}, \textbf{V}, the down-triangle \(\forall\)
(universal quantifier) and the horizontal arms \(\exists\)/\textsf{E} —
fitting, since Lean and its quantifiers do a lot of the work. It is a
deliberate construction on the unit circle and the golden ratio
(\(\varphi = 0.618\)), not a hand-traced drawing.

\begin{enumerate}[leftmargin=*,itemsep=1pt]
  \item Equilateral triangle on the unit circle, apex down \(\to T_{\mathrm{out}}\).
  \item Copy scaled to circumradius \(\varphi \to T_{\mathrm{in}}\); \(\mathrm{RING}=T_{\mathrm{out}}-T_{\mathrm{in}}\).
  \item Band thickness \(T=\tfrac12(1-\varphi)\) is the one gap/stroke unit.
  \item \textbf{Top-right cut} — the right-leg segment whose height equals its
        own base (\(2T/\sqrt3\)): a 60\(^\circ\) parallelogram on the triangle's
        inner \& outer "/". The flush line \(L\) derives from it.
  \item \textbf{Crossbar} — the F arm / \(\forall\) bar / \(\exists\) stroke, "/"-beveled.
  \item \textbf{Arrow cut} — a \(\perp\) cut through a chosen "\textbackslash" midpoint.
  \item \textbf{Accent + tertiary} fills complete the figure.
\end{enumerate}

% ── the family ───────────────────────────────────────────────────────
\section*{The family}
Both members share the construction; the accent and field colors invert.

\begin{center}
\begin{tabular}{>{\centering\arraybackslash}m{42mm} >{\centering\arraybackslash}m{42mm}}
  \includegraphics[width=38mm]{fastverk_full_1024.png} &
  \includegraphics[width=38mm]{fastverk_lower_1024.png} \\[2mm]
  \textbf{full} & \textbf{lower} \\
  {\footnotesize accent field + muted cut} & {\footnotesize muted field + accent arrow} \\
  {\footnotesize heraldic; large sizes} & {\footnotesize mark-like; small sizes} \\
\end{tabular}
\end{center}

% ── color ────────────────────────────────────────────────────────────
\section*{Color}
\fvswatch{fvbg}{15161A}\quad
\fvswatch{fvfg}{ECE7DA}\quad
\fvswatch{fvaccent}{E0A33E}\quad
\fvswatch{fvtertiary}{4A565A}

\medskip
\noindent\textbf{bg} background \quad \textbf{fg} mark \quad
\textbf{accent} arrow/field \quad \textbf{tertiary} muted cut/field.
The tertiary is deliberately desaturated so it recedes behind the accent.

% ── clear space & minimum size ───────────────────────────────────────
\section*{Clear space \& minimum size}
The rounded-rect background \emph{is} the clear space (corner radius
\(0.20\times\) the icon side). Both variants hold the silhouette down to
32\,px; below that, prefer \textbf{lower} (the arrow gives a clearer focal point).

\begin{center}
\begin{tabular}{cccc}
  \includegraphics[width=8mm]{fastverk_lower_16.png} &
  \includegraphics[width=12mm]{fastverk_lower_32.png} &
  \includegraphics[width=18mm]{fastverk_lower_64.png} &
  \includegraphics[width=26mm]{fastverk_lower_128.png} \\
  16 & 32 & 64 & 128 \\
\end{tabular}
\end{center}

% ── platform assets ──────────────────────────────────────────────────
\section*{Platform assets}
All built hermetically from the mark (\texttt{bazel build //icons:icon\_set}):

\begin{center}
\begin{tabular}{ll}
  \textbf{macOS} & \texttt{fastverk\_\{full,lower\}.icns} (16--1024, @2x) \\
  \textbf{Windows} & \texttt{fastverk\_\{full,lower\}.ico} (16--256) \\
  \textbf{Web / favicon} & \texttt{fastverk\_lower.ico}, \texttt{*.png} 16--1024 \\
  \textbf{Icon Composer} & layered SVG: \texttt{.bg/.tint/.arrow/.mark.svg} \\
\end{tabular}
\end{center}

\end{document}
